\documentclass[11pt,letterpaper]{article}
\usepackage[T1]{fontenc}
\usepackage{lmodern}
\usepackage[margin=1in]{geometry}
\usepackage{amsmath,amssymb,amsthm,mathtools}
\usepackage{booktabs,longtable,array}
\usepackage{microtype}
\usepackage{xcolor}
\usepackage{hyperref}
\usepackage{fancyhdr}
\usepackage{listings}
\usepackage{enumitem}
\definecolor{inkblue}{RGB}{27,57,88}
\definecolor{quietgray}{RGB}{85,85,85}
\hypersetup{colorlinks=true,linkcolor=inkblue,citecolor=inkblue,urlcolor=inkblue,
 pdftitle={The Exponential Growth Constant of OEIS A158415},
 pdfauthor={Research article}}
\pagestyle{fancy}
\fancyhf{}
\fancyhead[L]{\small\textsc{The growth constant of A158415}}
\fancyhead[R]{\small September 2026}
\fancyfoot[C]{\thepage}
\setlength{\headheight}{14pt}
\setlength{\parskip}{4pt}
\setlength{\parindent}{1.2em}
\setcounter{tocdepth}{2}
\numberwithin{equation}{section}
\newtheorem{theorem}{Theorem}[section]
\newtheorem{lemma}[theorem]{Lemma}
\newtheorem{proposition}[theorem]{Proposition}
\newtheorem{corollary}[theorem]{Corollary}
\theoremstyle{definition}
\newtheorem{definition}[theorem]{Definition}
\newtheorem{remark}[theorem]{Remark}
\newcommand{\Sset}{\mathcal S}
\newcommand{\Z}{\mathbb Z}
\newcommand{\Q}{\mathbb Q}
\newcommand{\N}{\mathbb N}
\newcommand{\OK}{\mathcal O_K}
\newcommand{\size}[1]{\lvert #1\rvert}
\newcommand{\coeff}[2]{[z^{#1}]#2}
\newcommand{\valuation}{v}
\DeclareMathOperator{\lcm}{lcm}
\DeclareMathOperator{\Norm}{N}
\lstset{basicstyle=\small\ttfamily,breaklines=true,columns=fullflexible,
 frame=single,rulecolor=\color{quietgray},xleftmargin=0.5em,xrightmargin=0.5em}

\title{\textbf{The Exponential Growth Constant\\of OEIS A158415}\\[0.7em]
\large Existence, an effective characterization,\\and certified numerical bounds}
\author{Research article}
\date{September 20, 2026}

\begin{document}
\maketitle
\begin{abstract}
Let $a(n)$ count the distinct real values represented by expressions with exactly
$n$ symbols, using the constant $1$, principal square root, and binary addition.
Numerical equality, rather than equality of expression trees, is essential.
We prove that a computable constant $\gamma>1$ exists such that
$\log a(n)=n\log\gamma+o(n)$, and obtain the certified bounds
$1.86204<\gamma<1.9208$.
The main algebraic step is a uniform injectivity theorem: any finite alphabet
of representable values can be used as digits in nested radicals after a
suitable fixed seed is chosen. Its proof uses an odd valuation at a prime
above $2$ and irreducible polynomials of controlled degree.
A complementary decomposition of minimum-size expressions yields convergent
finite-data upper bounds. Together these arguments also prove the exact
critical identity $\sum_{n\ge1}a(n)\gamma^{-n}=\gamma^2$; in particular,
$a(n)=o(\gamma^n)$.
The numerical lower bound is backed by an independently replayable certificate
for $1{,}526{,}536$ pairwise distinct values, using only integer arithmetic.
A separate combinatorial majorant certifies the upper bound and, after local
simplifications, proves the first 21 listed sequence terms exactly.
The value of $\gamma$ beyond the stated interval, a ratio limit, and a possible
$n^{-3/2}$ prefactor are not determined here.
\end{abstract}

\clearpage
{\small\setlength{\parskip}{0pt}
\tableofcontents}
\clearpage

\section{The sequence and the question}

An expression is built by the grammar
\[
 E ::= 1 \;\mid\; \sqrt E \;\mid\; (E+E).
\]
Its size is the number of occurrences of $1$, square-root operators, and
addition operators. Parentheses are not counted. Thus
\[
 |1|=1,\qquad |\sqrt E|=|E|+1,\qquad |E+F|=|E|+|F|+1.
\]
Every square root denotes the nonnegative real square root. All values produced
by this grammar are at least $1$.
The sequence in the question is OEIS A158415 \cite{oeis}, introduced by
Vladimir Reshetnikov in 2009. We study its defining counting problem, not just
an extrapolation of its recorded initial values.

Write $\Sset$ for the set of all represented values and define
\[
 c(x)=\min\{|E|:E\text{ represents }x\},\qquad x\in\Sset.
\]
Replacing a leaf $1$ by $\sqrt1$ increases the size by one without changing
the value. Consequently, the exact-size and cumulative interpretations agree:
\begin{equation}\label{eq:cumulative}
 a(n)=\#\{x\in\Sset:c(x)\le n\}.
\end{equation}
Set $a(0)=0$ and
\begin{equation}\label{eq:bdef}
 b(n)=a(n)-a(n-1)=\#\{x\in\Sset:c(x)=n\}.
\end{equation}

The first few values are easy to check. At size at most four the values are
$1$, $2$, and $\sqrt2$, so $a(4)=3$. The main difficulty is that different
syntactic expressions can have the same value. Besides associativity and
commutativity, one has identities such as
\[
 \sqrt{1+1+1+1}=1+1,
 \qquad
 \sqrt{3+2\sqrt2}=1+\sqrt2.
\]
Here integer coefficients and integer constants are abbreviations for sums
of copies of $1$; they are not additional symbols in the grammar.

The original sequence discussion explicitly raised the issue of proving
numerical equality and questioned the fixed tolerance used in a PARI
calculation \cite{discussion}. That distinction remains important here:
our numerical certificates never infer equality merely from overlapping
approximations.

\section{Main result and its scope}

\begin{theorem}[Growth constant and effective characterization]\label{thm:main}
There is a computable real number $\gamma$ such that
\begin{equation}\label{eq:mainlimit}
 \lim_{n\to\infty}a(n)^{1/n}=\gamma,
 \qquad
 a(n)=\gamma^{n+o(n)}.
\end{equation}
It satisfies
\begin{equation}\label{eq:numericalbounds}
 1.86204<\gamma<1.9208.
\end{equation}
For each $m\ge1$, let $r_m\in(0,1)$ be the unique solution of
\begin{equation}\label{eq:rm}
 r_m\left(1+\sum_{j=1}^m b(j)r_m^{j+1}\right)=1,
 \qquad \lambda_m=r_m^{-1}.
\end{equation}
Then $\lambda_m$ is nondecreasing and
\begin{equation}\label{eq:sup}
 \gamma=\sup_{m\ge1}\lambda_m=\lim_{m\to\infty}\lambda_m.
\end{equation}
There are also explicitly defined algebraic upper bounds $\mu_h\ge\gamma$,
computed from $b(1),\ldots,b(3h-1)$, with $\mu_h\to\gamma$.
Finally,
\begin{equation}\label{eq:criticalsum}
 \boxed{\displaystyle \sum_{n\ge1}a(n)\gamma^{-n}=\gamma^2,}
 \qquad a(n)=o(\gamma^n).
\end{equation}
\end{theorem}

The finite lower-bound equation can equivalently be written as the polynomial
identity
\begin{equation}\label{eq:lambda-polynomial}
 \lambda_m^{m+2}=\lambda_m^{m+1}
       +\sum_{j=1}^m b(j)\lambda_m^{m-j}.
\end{equation}
Thus the theorem gives more than a statement that the sequence lies between
two exponential functions. It proves a single exponential rate, supplies
convergent finite-data characterizations from both sides, and identifies an
exact normalization of the generating function at its radius of convergence.

It does \emph{not} give a closed form for $\gamma$, prove a limit for
$a(n+1)/a(n)$, or establish an equivalent of the form
$C\gamma^n n^{-3/2}$. These stronger assertions must not be inferred from the
root limit. In fact, \eqref{eq:criticalsum} excludes a positive constant
prefactor $C\gamma^n$. More generally, if
$a(n)\sim C\gamma^n n^{-\theta}$ with $C>0$, then necessarily $\theta>1$.

\section{Elementary structural facts}

\subsection{Size, magnitude, and algebraic integrality}

\begin{lemma}\label{lem:basic}
Every represented value is an algebraic integer. If an expression has $L$
leaves and $U$ square-root nodes, then its size is $2L-1+U$ and every conjugate
of its value has absolute value at most $L$.
\end{lemma}
\begin{proof}
The number of binary nodes is $L-1$, giving the size formula. Algebraic
integers are closed under addition, and a square root of an algebraic integer
is integral over the ring of algebraic integers and hence over $\Z$.
For the conjugate bound, put all subexpression values in a finite number field
and apply any embedding into $\mathbb C$. Addition obeys the triangle
inequality. For a radical node,
$|\sigma(\sqrt x)|=\sqrt{|\sigma(x)|}\le\sqrt L\le L$.
The assertion follows by induction from the leaf value $1$.
\end{proof}

In particular, $c(k)=2k-1$ for each positive integer $k$: at least $k$ leaves
are required to attain magnitude $k$, and their sum realizes this size.
There are therefore new values at arbitrarily large sizes.

A prefix encoding of an expression is a word over three symbols. This gives
the coarse but sufficient initial estimate
\begin{equation}\label{eq:coarse}
 1\le a(n)\le 3^n.
\end{equation}

\subsection{Multiplication is available at finite cost}

Multiplication is not an allowed symbol. Nevertheless its values are available.
This observation is needed for the number-field argument, not as a claim that
multiplication is inexpensive.

\begin{lemma}\label{lem:mult}
If $x,y\in\Sset$, then $xy\in\Sset$.
\end{lemma}
\begin{proof}
Induct on the sum of the sizes of chosen expressions for $x$ and $y$.
A factor $1$ is harmless. If either expression is a sum, distribute the product
and use the induction hypothesis. Otherwise both nontrivial factors are
radicals, and
\[
 (\sqrt u)(\sqrt v)=\sqrt{uv}.
\]
The product $uv$ involves smaller expressions, so it too is representable.
\end{proof}

It follows that every element of a ring $\Z[d_1,\ldots,d_s]$ with
$d_i\in\Sset$ can be written as $A-B$, where $A,B\in\Sset\cup\{0\}$.
If desired, adding $1$ to both parts makes $A$ and $B$ genuine elements of
$\Sset$.

\section{Injective coding by nested radicals}

This section is the central algebraic argument. It overcomes the fact that
ordinary addition is not an injective pairing operation.

We use standard facts about orders in number fields and discrete valuations:
if $K$ is a number field and $\mathfrak p$ lies above $2$, then the normalized
valuation $v=v_{\mathfrak p}$ takes integer values on $K^\times$, its residue
field is finite of characteristic $2$, and it extends to algebraic extensions.
An order in $K$ has finite index in $\OK$. These facts, and the corresponding
Newton-polygon viewpoint, are treated in \cite[Chapters 2, 3, and 7]{milne}.
The particular irreducibility argument below is proved directly.

\subsection{A representable element with odd valuation}

\begin{lemma}\label{lem:odd}
Let $D\subset\Sset$ be finite and $K=\Q(D)$. There are a prime
$\mathfrak p$ of $K$ above $2$ and an element $\pi\in\Sset\cap K$ such that
$h=v_{\mathfrak p}(\pi)$ is a positive odd integer.
\end{lemma}
\begin{proof}
Choose any prime $\mathfrak p$ above $2$, and write $e=v_{\mathfrak p}(2)$.
If $e$ is odd, take $\pi=2$.

Suppose $e$ is even. The ring $R=\Z[D]$ is an order in $K$: it is generated
by algebraic integers and spans $K$ over $\Q$. Choose a positive integer $I$
with $I\OK\subset R$, and choose $u\in\OK$ with $v(u)=1$.
Then
\[
 t=Iu\in R,\qquad v(t)=e\,v_2(I)+1=:h,
\]
so $h$ is odd. By Lemma~\ref{lem:mult}, write $t=A-B$ with
$A,B\in\Sset$. For a sufficiently large integer $M$,
$v(2^M B)>h$. The positive, representable number
\[
 \pi=A+(2^M-1)B=t+2^M B
\]
then has valuation $h$. Integer multiples here are realized by repeated
addition, so the construction stays inside the original language.
\end{proof}

The element $\pi$ need not be a uniformizer. Allowing an arbitrary positive
\emph{odd} valuation is what makes the construction work without requiring
the given order to be integrally closed at $2$.

\subsection{A local irreducibility lemma}

\begin{lemma}\label{lem:irreducible}
Let $K$ be as above, let $v(\pi)=h>0$ be odd, and choose a power of two
$N=2^t>h$. Suppose $Q(X)\in\OK[X]$ is monic of degree $d=2^r$ and
its reduction modulo $\mathfrak p$ has the form $X^d+c$. Then
\begin{equation}\label{eq:irreducible}
 Q(X)^N-\pi
\end{equation}
is irreducible over $K$.
\end{lemma}
\begin{proof}
The residue field is finite, so the map $x\mapsto x^d$ is bijective.
Choose a residue-field root of $X^d+c$ and lift it to $a\in\OK$.
For $q(X)=Q(X+a)$, reduction modulo $\mathfrak p$ is exactly $X^d$.
Thus $q$ is monic and every nonleading coefficient has valuation at least one.

Extend $v$ to an algebraic closure and let $\alpha$ be any root of
$q(X)^N-\pi$. Then $v(q(\alpha))=h/N\in(0,1)$.
If $v(\alpha)<0$, the leading term of $q(\alpha)$ has uniquely smallest
valuation, giving a negative valuation, a contradiction.
If $v(\alpha)\ge1/d$, every term of $q(\alpha)$ has valuation at least one,
also a contradiction. For $0\le v(\alpha)<1/d$, the leading term again has
uniquely smallest valuation, and therefore
\begin{equation}\label{eq:rootvaluation}
 v(\alpha)=\frac{h}{Nd}.
\end{equation}
In particular $v(\alpha)>0$.

If the polynomial had a monic factor of degree $k$ over $K$, the constant
coefficient of that factor would be, up to sign, a product of $k$ roots.
Its valuation would be $kh/(Nd)$, which must be an integer.
Since $h$ is odd and $Nd$ is a power of two, this implies $Nd\mid k$.
No nonconstant proper factor is possible. Translating back proves the claim.
\end{proof}

\subsection{One seed works for the whole finite alphabet}

\begin{theorem}[Uniform injectivity of radical spines]\label{thm:spines}
For every finite set $D\subset\Sset$ there is a representable positive seed
$s$ such that all the numbers
\begin{equation}\label{eq:spine}
 d_0+\sqrt{d_1+\sqrt{\cdots+\sqrt{d_r+s}}},
 \qquad r\ge0,\quad d_i\in D\cup\{0\},
\end{equation}
are distinct as the complete digit string $(d_0,\ldots,d_r)$ varies.
For $r=0$ the expression means $d_0+s$. A zero digit means that the
corresponding addition is omitted.
\end{theorem}
\begin{proof}
Take $K=\Q(D)$, choose $\pi$ by Lemma~\ref{lem:odd}, and put
$s=\pi^{1/N}$, where $N=2^t>v(\pi)$. This is a fixed representable seed,
obtained by $t$ successive square roots.

For the string in \eqref{eq:spine}, define
\[
 Q_0(X)=X-d_0,\qquad Q_j(X)=Q_{j-1}(X)^2-d_j\quad(1\le j\le r).
\]
The resulting value $x$ satisfies $Q_r(x)=s$, so it is a root of
\[
 P_{\boldsymbol d}(X)=Q_r(X)^N-\pi.
\]
Modulo $\mathfrak p$, repeated squaring gives a polynomial of the form
$X^{2^r}+c$. Lemma~\ref{lem:irreducible} shows that $P_{\boldsymbol d}$ is
monic irreducible of degree $N2^r$.

Equal values must have equal monic minimal polynomials over $K$.
Their degrees first force the same $r$. Equality of the minimal polynomials
then gives $Q_r^N=\widetilde Q_r^N$, hence $Q_r=\widetilde Q_r$, since the
polynomials are monic.
The decomposition $Q_j=Q_{j-1}^2-d_j$ is unique: if
$A^2-u=B^2-v$ with $A,B$ monic of the same positive degree, then
$(A-B)(A+B)=u-v$ is constant. Unless $A=B$, its left side has positive degree.
Thus $A=B$ and $u=v$.
Peeling off the successive squares recovers all digits, including $d_0$.
\end{proof}

\begin{remark}
The seed can have a very large expression size. This is harmless for an
asymptotic lower bound with a fixed finite alphabet, but it prevents the proof
from being interpreted as a practical short encoding of arbitrary finite
sets. The order of limits is important: first fix the alphabet and let the
output size tend to infinity; only then enlarge the alphabet.
\end{remark}

\section{Existence of the exponential rate}

\subsection{Finite alphabets give genuine exponential lower bounds}

Assign one expression of size $\ell(d)$ to each digit in a finite alphabet
$D$. Adding a nonzero digit costs $\ell(d)+1$ symbols, whereas the zero digit
costs nothing. Define
\begin{equation}\label{eq:digitgf}
 D_\ell(z)=1+\sum_{d\in D}z^{\ell(d)+1}.
\end{equation}
If the seed of Theorem~\ref{thm:spines} has an expression of size $L$, then
all its spines have the exact generating function
\begin{equation}\label{eq:spinegf}
 z^L\sum_{r\ge0}z^r D_\ell(z)^{r+1}
   =\frac{z^L D_\ell(z)}{1-zD_\ell(z)}.
\end{equation}
This counts \emph{distinct values}, not merely formal strings, by the
injectivity theorem. Summing over expression sizes at most $n$ introduces
an additional factor $(1-z)^{-1}$.

Let $\rho\in(0,1)$ solve $\rho D_\ell(\rho)=1$.
The zero of $1-zD_\ell(z)$ at $\rho$ is simple, and it is the only zero
on its circle of convergence. Indeed equality in the triangle inequality
would require the degree-one term of $zD_\ell(z)$ to have positive real phase,
forcing $z=\rho$. Partial fractions, or the usual simple-pole coefficient
calculation, therefore gives
\begin{equation}\label{eq:finitealphabetlower}
 a(n)\ge C_D\rho^{-n}\qquad(n\text{ sufficiently large})
\end{equation}
for a constant $C_D>0$. The standard rational-function framework is explained
in \cite[Chapters IV--V]{fs}; here the needed pole properties are explicit.

Taking all values with $c(d)\le m$, and assigning their minimum sizes,
gives the lower bases $\lambda_m$ in \eqref{eq:rm}.
In particular,
\begin{equation}\label{eq:liminflambda}
 \liminf_{n\to\infty}a(n)^{1/n}\ge \sup_m\lambda_m.
\end{equation}

\subsection{The generating-function argument that removes oscillation}

Define
\begin{equation}\label{eq:AB}
 A(z)=\sum_{n\ge1}a(n)z^n,\qquad
 B(z)=\sum_{n\ge1}b(n)z^{n+1}=z(1-z)A(z).
\end{equation}
Put $g=\sup_m\lambda_m$. It is finite by \eqref{eq:coarse} and
\eqref{eq:liminflambda}, and $g>1$ already follows from the alphabet $\{1\}$.
For every $0<z\le1/g$ and every $m$,
\[
 z\left(1+\sum_{j=1}^m b(j)z^{j+1}\right)\le1.
\]
Monotone convergence in $m$ shows that $B(z)$ is finite and
\begin{equation}\label{eq:boundaryinequality}
 z(1+B(z))\le1 \qquad(0<z\le1/g).
\end{equation}
Thus $A(z)$ is finite there as well. By the elementary radius formula for a
power series with nonnegative coefficients,
\[
 \limsup_{n\to\infty}a(n)^{1/n}\le g.
\]
Together with \eqref{eq:liminflambda}, this proves the root limit, with
$\gamma=g$. Writing $R=1/\gamma$, we have established
\begin{equation}\label{eq:initial-boundary}
 B(R)<\infty,\qquad R(1+B(R))\le1,
 \qquad A(R)\le\gamma^2.
\end{equation}
The inequality will become an equality in Section~\ref{sec:effective}.
Already it implies $a(n)\gamma^{-n}\to0$.

\subsection{An elementary explicit subfamily}

For integer digits, a single seed works for the entire infinite alphabet:
take $K=\Q$, $\pi=2$, $N=2$, and $s=\sqrt2$, of size four.
The proof of Theorem~\ref{thm:spines} applies uniformly to all nonnegative
integer digits. A digit $k$ costs $2k$ symbols when added to a spine, so
$D(z)=(1-z^2)^{-1}$. Formula~\eqref{eq:spinegf} becomes
\[
 \frac{z^4}{1-z-z^2}.
\]
With $F_1=F_2=1$, cumulative counting gives the entirely explicit estimate
\begin{equation}\label{eq:fibonacci}
 a(n)\ge F_{n-1}-1\qquad(n\ge4).
\end{equation}
In particular, $\gamma\ge(1+\sqrt5)/2$. The numerical certificate below
improves this substantially.

\section{Finite-data upper bounds and computability}\label{sec:effective}

The matching upper mechanism uses addition's associativity. It is here that
counting numerical values, with minimum expression sizes, supplies useful
structure rather than just an obstacle.

\subsection{Additive weights and minimum subexpressions}

Give an expression the weight $w(E)=|E|+1$, and a value the minimum weight
$w(x)=c(x)+1$. Then
\[
 w(E+F)=w(E)+w(F),\qquad w(\sqrt E)=w(E)+1.
\]
Flatten every maximal addition into a multiset of atoms, where an atom is
$1$ or a square-root expression.
If the entire expression has minimum size, every subexpression is minimum.
More strongly, the sum of \emph{any nonempty submultiset} of a flattened
addition is minimum for its value. Otherwise one could regroup that
submultiset, replace it by a smaller expression, and shorten the whole
expression. Regrouping does not change the number of addition symbols.

For an integer $h\ge2$, define the finite polynomials
\begin{align}
 \mathcal D_h(z)&=1+\sum_{n+1\le h}b(n)z^{n+1},\label{eq:Dh}\\
 \mathcal T_h(z)&=\sum_{h<n+1\le3h}b(n)z^{n+1}.\label{eq:Th}
\end{align}

\begin{lemma}[Large-expression decomposition]\label{lem:decompose}
Every minimum-size expression whose weight is greater than $3h$ can be
decomposed, without changing its weight, in one of two ways:
\begin{enumerate}[label=(\roman*),itemsep=2pt]
\item as $d+\sqrt E$, where $d=0$ or $w(d)\le h$, and $w(E)>h$;
\item as $E+F$, where both $w(E)$ and $w(F)$ exceed $h$.
\end{enumerate}
All the nonzero components may be represented with their minimum weights.
\end{lemma}
\begin{proof}
A radical at the root gives (i) with $d=0$.
Otherwise flatten the top addition. If its largest atom has weight greater
than $h$, consider the total weight of the remaining atoms. If that remainder
exceeds $h$, split it off to obtain (ii). If it is at most $h$, the large atom
cannot be $1$ and is therefore a radical, giving (i).

If every atom has weight at most $h$, collect atoms until their total weight
first exceeds $h$. This total is at most $2h$, and the remaining atoms have
weight greater than $h$ because the whole weight exceeds $3h$. This gives
(ii). Minimum-weight status follows from the preceding submultiset argument.
In case (i), the radical child has weight greater than $2h-1$, hence greater
than $h$.
\end{proof}

Stop the decomposition when a component has weight in $(h,3h]$.
The resulting codes are counted, with possible overcounting, by the
nonnegative formal power series $H_h$ satisfying
\begin{equation}\label{eq:H}
 H_h=\mathcal T_h+z\mathcal D_hH_h+H_h^2,
 \qquad H_h(0)=0.
\end{equation}
The last term uses ordered children; allowing both orders only enlarges the
upper count. Each value of weight greater than $h$ has at least one code of
its minimum weight, and a code determines a value. Therefore
\begin{equation}\label{eq:Hmajorant}
 \sum_{n+1>h}b(n)z^{n+1}\preceq H_h(z),
\end{equation}
where $\preceq$ denotes coefficientwise domination.

\subsection{Explicit algebraic upper bounds}

Let $s_h$ be the first positive solution, before
$z\mathcal D_h(z)=1$, of
\begin{equation}\label{eq:sh}
 (1-z\mathcal D_h(z))^2=4\mathcal T_h(z),
 \qquad \mu_h=s_h^{-1}.
\end{equation}
Such a solution is unique: before $z\mathcal D_h(z)=1$ the left side decreases
strictly and the right side increases. Also $\mathcal T_h$ is nonzero, since
positive integers have minimum weights $2,4,6,\ldots$.
The smaller solution of \eqref{eq:H} is
\begin{equation}\label{eq:Hsolution}
 H_h(z)=\frac{1-z\mathcal D_h(z)
 -\sqrt{(1-z\mathcal D_h(z))^2-4\mathcal T_h(z)}}2.
\end{equation}
It is finite at $s_h$ and has nonnegative coefficients. For example, bounded
iteration of \eqref{eq:H} proves convergence for positive $z\le s_h$ without
any assumptions about other complex singularities.
Equation~\eqref{eq:Hmajorant} gives
\begin{equation}\label{eq:mubound}
 \gamma\le\mu_h.
\end{equation}

\begin{proposition}\label{prop:muconverges}
The upper bounds satisfy $\mu_h\to\gamma$. Consequently
\[
 \gamma=\sup_m\lambda_m=\inf_{h\ge2}\mu_h.
\]
The running minima of $\mu_h$ form a nonincreasing sequence of upper bounds.
\end{proposition}
\begin{proof}
Fix $0<z<R=1/\gamma$. By \eqref{eq:initial-boundary} and strict monotonicity,
$z(1+B(z))<1$. As $h\to\infty$,
\[
 \mathcal D_h(z)\longrightarrow1+B(z),\qquad
 \mathcal T_h(z)\longrightarrow0.
\]
Thus $(1-z\mathcal D_h(z))^2>4\mathcal T_h(z)$ for all sufficiently large
$h$, and $s_h>z$. Hence $\limsup_h\mu_h\le1/z$. Letting $z\uparrow R$ and
using \eqref{eq:mubound} proves the result.
\end{proof}

\subsection{The exact critical identity}

\begin{proposition}\label{prop:critical}
At $R=1/\gamma$,
\begin{equation}\label{eq:criticalD}
 R(1+B(R))=1.
\end{equation}
Equivalently, $A(R)=R^{-2}$.
\end{proposition}
\begin{proof}
We already know the left side is at most one and $B(R)$ is finite.
Suppose it were strictly less than one. Then, as $h\to\infty$,
$\mathcal T_h(R)\to0$ and $1-R\mathcal D_h(R)$ stays bounded away from zero.
For some finite $h$ we would have
\[
 1-R\mathcal D_h(R)>0,\qquad
 (1-R\mathcal D_h(R))^2>4\mathcal T_h(R).
\]
Both inequalities are strict polynomial inequalities. They would remain true
at some $z>R$, implying convergence of $H_h(z)$ and hence of $B(z)$.
This contradicts the definition of $R$ as the radius of convergence.
Therefore \eqref{eq:criticalD} holds. Substitution of
$B(R)=R(1-R)A(R)$ gives $A(R)=R^{-2}$.
\end{proof}

\subsection{Why this is an effective characterization}

Every finite set of expression values is exactly computable using algebraic
number arithmetic; a completely explicit separation bound is proved in
Section~\ref{sec:separation}. Thus all $b(1),\ldots,b(M)$ are computable in
finite time. Positive roots of the finite polynomials
\eqref{eq:rm} and \eqref{eq:sh} can be bracketed by rational arithmetic.

To approximate $\gamma$ within a prescribed positive error $\varepsilon$,
compute increasing finite sets of exact terms and the associated lower bounds
$\lambda_m$ and running upper bounds $\min_{2\le h\le H}\mu_h$.
Stop when the two rational brackets differ by less than $\varepsilon$.
Proposition~\ref{prop:muconverges} proves that this algorithm terminates.
No unproved convergence rate or numerical extrapolation is needed.
This is a computability statement, not a polynomial-time claim.

\section{A certified numerical lower bound}

\subsection{What the certificate proves}

The archive contains a directed acyclic graph of $1{,}526{,}536$ expressions,
each with an explicitly recorded symbol cost at most 28. Reusing a stored
node is only a compression technique: the cost calculation charges a child
again every time it is used, so this is \emph{formula size}, not circuit size.

For each node the independent verifier computes an interval
$[\ell/2^{256},u/2^{256}]$ using integer arithmetic. For addition it adds
endpoints. For a square root it uses
\begin{equation}\label{eq:intervalsqrt}
 \ell'=\left\lfloor\sqrt{\ell\,2^{256}}\right\rfloor,
 \qquad
 u'=\left\lceil\sqrt{u\,2^{256}}\right\rceil.
\end{equation}
An integer square-root routine computes these bounds exactly.
After sorting all intervals, the verifier checks that every consecutive pair
is strictly disjoint. This proves pairwise distinctness of all stored values.
The smallest certified gap in this sample is approximately
$5.46212307648\times10^{-13}$; the exact rational gap is recorded in the
verification output.

The generator discards an interval whenever it overlaps one already retained.
Such a discard is conservative: it is \emph{not} an assertion that the two
values are equal. Therefore the output is unconditionally a lower-bound
sample, even if a very close distinct value were discarded.

Let $q_n$ be the number of sample nodes assigned cost $n$.
They need not be assumed to have minimum cost. Theorem~\ref{thm:spines}
can use these assigned costs directly, so the positive root of
\begin{equation}\label{eq:samplepoly}
 z\left(1+\sum_{n=1}^{28}q_nz^{n+1}\right)=1
\end{equation}
gives a valid lower bound on $\gamma$. In particular, no completeness claim
about the enumeration is required.

\subsection{The finite-root calculation}

The cumulative sample counts agree with the 28 values currently listed in
the OEIS entry \cite{oeis}. For the lower-bound proof, however, the stored
expressions and interval verification replace reliance on those listed
numbers.
The reciprocal root of \eqref{eq:samplepoly} is bracketed by exact rational
bisection and begins
\begin{equation}\label{eq:samplelambda}
 1.86204402268771032792879\ldots.
\end{equation}
At the exact rational number $z=25000/46551=1/1.86204$, the left side of
\eqref{eq:samplepoly}, minus one, is strictly positive; its value is
approximately $1.09151356699874\times10^{-5}$. Since that left side is
strictly increasing in $z$, this proves $\gamma>1.86204$.

\begin{table}[ht]
\centering
\begin{tabular}{rr}
\toprule
Maximum assigned digit cost & Reciprocal root, displayed approximately\\
\midrule
4  & 1.618033988749895\\
8  & 1.741298211719966\\
12 & 1.792996772315302\\
16 & 1.821892848098540\\
20 & 1.840189044447032\\
24 & 1.852813818187051\\
28 & 1.862044022687710\\
\bottomrule
\end{tabular}
\caption{Certified finite-alphabet lower bases. The archive records rational
brackets, not just these decimal displays.}
\label{tab:lower}
\end{table}

The constants in the accompanying eventual inequalities
$a(n)\ge C_m\lambda_m^n$ are not computed: the general seed construction can
be large. The certified conclusion concerns the limiting exponential rate.

\section{A certified upper bound below 1.921}

\subsection{An unordered-sum majorant}

Consider a class of formal expressions in which maximal additions are
flattened into multisets and a radical of a pure integer square is forbidden.
Every minimum-size expression belongs to this class: replacing
$\sqrt{q^2}$ by $q$ is always shorter. Other numerical identities are ignored,
so the class provides an upper bound.

Use weight $|E|+1$, and let $U(z)=\sum_{j\ge2}u_jz^j$ count these formal
nonempty expressions. The atom $1$ has weight two. A radical adds one to the
weight of its child, except that the forbidden radicals contribute
\begin{equation}\label{eq:P}
 P(z)=\sum_{q\ge1}z^{2q^2+1}.
\end{equation}
Consequently the atom series is $z^2+zU(z)-P(z)$, and the multiset construction
gives
\begin{equation}\label{eq:Ugf}
 1+U(z)=\exp\left(\sum_{k\ge1}
   \frac{z^{2k}+z^kU(z^k)-P(z^k)}k\right).
\end{equation}
The formula can also be checked directly by multiplying
$(1-z^j)^{-p_j}$ over atom weights $j$, where $p_j$ is the number of atoms
of that weight; see \cite[Chapter I]{fs} for the general multiset construction.

Minimum-size expressions inject into this formal class after one
representative per value is selected. Therefore
\begin{equation}\label{eq:au}
 a(n)\le\sum_{j=2}^{n+1}u_j.
\end{equation}

\subsection{Exact coefficients by the Euler transform}

Let $c_0=1$ and $c_n=u_n$ for $n\ge1$. The atom counts are
\[
 p_n=\boldsymbol1_{n=2}+u_{n-1}
       -\#\{q\ge1:n=2q^2+1\}.
\]
Put $e_k=\sum_{d\mid k}d\,p_d$. Logarithmic differentiation of the product
formula gives the integer recurrence
\begin{equation}\label{eq:euler}
 n c_n=\sum_{k=1}^n e_kc_{n-k}.
\end{equation}
There is no circular dependence: $p_n$ uses only $u_{n-1}$.
The certificate computes these coefficients through $n=200$, checking
nonnegativity and exact divisibility at every step.

\subsection{A convergence criterion proved by radical height}

For $r>0$ with $r^2<1/3$, define
\begin{align}
 \mathcal H(r)&=\sum_{k\ge2}
   \frac{r^{2k}+r^kU(r^k)-P(r^k)}k,\label{eq:Htaildef}\\
 \mathcal F(r)&=\log r+1+r^2-r-P(r)+\mathcal H(r).\label{eq:Fcriterion}
\end{align}
All evaluations $U(r^k)$ in this definition converge: the number of formal
objects of weight $j$ is at most $3^{j-1}$, by prefix encoding of an associated
ordered expression. Hence $U(x)$ converges for $x<1/3$.

\begin{proposition}\label{prop:uppercriterion}
If
\begin{equation}\label{eq:uppercriterion}
 r^2+r<1,\qquad r^2<1/3,\qquad \mathcal F(r)<0,
\end{equation}
then $U(r)$ is finite and $\gamma<1/r$.
\end{proposition}
\begin{proof}
Let $U_j$ count objects of radical height at most $j$.
At height zero they are positive integers, so
$1+U_0(r)=(1-r^2)^{-1}<1/r$.
Every pure integer square is already present at height zero; thus removing
$P$ in each later atom construction is legitimate and leaves nonnegative
coefficients.
The classes increase with height, and their union is the whole formal class.

Writing $C_j=1+U_j$, the $k=1$ term of the multiset recurrence gives
\[
 C_{j+1}(r)\le
 \exp\bigl(rC_j(r)+r^2-r-P(r)+\mathcal H(r)\bigr).
\]
For $k\ge2$, the bounded-height classes were bounded by the full class at
$r^k$, where convergence is already known. If $C_j(r)\le1/r$, the exponent
is at most $1+r^2-r-P(r)+\mathcal H(r)<-\log r$.
Induction gives $C_j(r)<1/r$ for every $j$, uniformly. Taking the increasing
limit proves convergence at $r$.

The estimates are strict and persist at a slightly larger positive argument:
the small-argument series in \eqref{eq:Htaildef} are locally uniformly
convergent while $r^2<1/3$. The same height argument thus proves convergence
at some $r'>r$. Equation~\eqref{eq:au} yields $\gamma\le1/r'<1/r$.
\end{proof}

\subsection{Explicit bounds for every omitted tail}

For $0<x<1/3$, the prefix bound gives
\begin{equation}\label{eq:Utail}
 U(x)\le\sum_{j=2}^{N}u_jx^j
       +\frac{3^N x^{N+1}}{1-3x}.
\end{equation}
Use the finite lower bound
$P_Q(x)=\sum_{q=1}^Qx^{2q^2+1}\le P(x)$ wherever $-P(x)$ occurs.
For the terms with $k>K$ in \eqref{eq:Htaildef}, omit the negative terms and
use $U(x)\le3x^2/(1-3x)$. Their total is at most
\begin{equation}\label{eq:Htailbound}
 \frac{r^{2(K+1)}}{(K+1)(1-r^2)}
 +\frac{3r^{3(K+1)}}{(K+1)(1-r^3)(1-3r^{K+1})}.
\end{equation}
Finally, for $t=(1-r)/(1+r)\in(0,1)$,
\[
 \log r=-2\sum_{j\ge0}\frac{t^{2j+1}}{2j+1}.
\]
After $J$ terms the positive sum's omitted tail is at most
\begin{equation}\label{eq:logtail}
 \frac{t^{2J+1}}{(2J+1)(1-t^2)}.
\end{equation}
These formulas reduce the entire verification to finitely many rational
operations and an explicitly bounded series for the logarithm.

\subsection{The numerical certificate}

The program \texttt{certify\_upper\_bound.py} evaluates an upper bound for
$\mathcal F(r)$ at the exact rational argument
\[
 r=\frac{1250}{2401}=\frac1{1.9208}
\]
using $N=200$, $K=40$, $Q=8$, and $J=100$.
All arithmetic is enclosed by outward-rounded dyadic intervals with
1024 fractional bits. Every tail is bounded by the displayed formulas.
The resulting upper bound is strictly negative and is approximately
\begin{equation}\label{eq:Fvalue}
 -8.442263650713085\times10^{-6}.
\end{equation}
The initial height margin $1/r-(1-r^2)^{-1}$ is positive, approximately
$0.5489798948$.
The exact signed integer numerator of the upper endpoint is stored in the
archive. Thus Proposition~\ref{prop:uppercriterion}, rather than an unverified
floating-point root finder, proves
\[
 \boxed{\gamma<1.9208.}
\]
Together with Section 7 this completes the numerical part of
Theorem~\ref{thm:main}.

\section{Sharper finite-term bounds and exact initial terms}

Two further local simplifications significantly tighten the formal upper
counts. They are useful for certifying initial terms even though the simpler
majorant of Section 8 suffices for the stated growth bound.

\subsection{Square factors in additive multiplicities}

If every multiplicity in a radical's flattened radicand is divisible by
$q^2$, $q\ge2$, then its value is $\sqrt{q^2t}=q\sqrt t$.
If $w(t)$ denotes the syntactic weight of the divided multiset, the old and
new weights are respectively
\[
 q^2w(t)+1,\qquad q(w(t)+1).
\]
The old weight exceeds the new one by $(q-1)(qw(t)-1)>0$.
No minimum-size expression can contain such an atom.

The allowed radicands therefore have squarefree gcd of their multiplicities.
If $V(z)$ counts the resulting formal expressions, the radicand series with
this restriction is
\begin{equation}\label{eq:mobius}
 \sum_{q\ge1}\mu(q)V(z^{q^2}),
\end{equation}
where $\mu$ is the ordinary Moebius function. Indeed a multiset with
multiplicity gcd $g$ is counted with coefficient
$\sum_{q^2\mid g}\mu(q)$, which is one precisely when $g$ is squarefree.
The coefficients of \eqref{eq:mobius} are consequently nonnegative despite
its alternating presentation.

\subsection{A family of denesting identities}

Put $\alpha_k=2^{1/2^k}$ for $k\ge1$. Its displayed expression has size $k+3$.
For $m\ge1$, forbid the atom obtained from the expanded square
\begin{equation}\label{eq:family}
 \sqrt{m^2+2m\alpha_k+\alpha_k^2}=m+\alpha_k.
\end{equation}
The right side is strictly shorter.
The atom on the left has weight
\begin{equation}\label{eq:familyweight}
 2m^2+8m+4+(2m+1)k.
\end{equation}
For $k=1$, combine $\alpha_1^2=2$ with the integer part; the multiplicity gcd
is one or two. For $k\ge2$, the distinct atom $\alpha_{k-1}$ has multiplicity
one. Thus none of these exclusions overlaps the square-factor restriction.
Their children are allowed dyadic-root atoms, and distinct pairs $(m,k)$
give distinct formal radicands.

Define
\begin{equation}\label{eq:Efamily}
 E(z)=\sum_{m\ge1}\sum_{k\ge1}
 z^{2m^2+8m+4+(2m+1)k}.
\end{equation}
The improved atom series is
\begin{equation}\label{eq:improvedatom}
 J(z)=z^2+z\sum_{q\ge1}\mu(q)V(z^{q^2})-z^3-E(z),
\end{equation}
and the expression series satisfies
\begin{equation}\label{eq:Vgf}
 1+V(z)=\exp\left(\sum_{j\ge1}\frac{J(z^j)}j\right).
\end{equation}
The term $-z^3$ removes $\sqrt1$. At every fixed degree all sums are finite.
The same integer Euler-transform recurrence computes the coefficients.
Every minimum-size expression avoids all these forbidden atoms, so $V$ is
still a valid coefficientwise upper model for minimum-size value counts.

\subsection{What is now proved about the listed terms}

Let $L_n$ denote the cumulative number of certified sample values and let
$U_n^*$ be the cumulative upper count from \eqref{eq:Vgf}. Then
$L_n\le a(n)\le U_n^*$. Equality holds between these bounds for every
$n\le21$. The remaining certified intervals are displayed below.

\begin{table}[ht]
\centering
\begin{tabular}{rrrr}
\toprule
$n$ & Certified lower bound $L_n$ & Certified upper bound $U_n^*$ & Gap\\
\midrule
16 & 1,379 & 1,379 & 0\\
17 & 2,424 & 2,424 & 0\\
18 & 4,277 & 4,277 & 0\\
19 & 7,588 & 7,588 & 0\\
20 & 13,513 & 13,513 & 0\\
21 & 24,162 & 24,162 & 0\\
22 & 43,336 & 43,337 & 1\\
23 & 77,978 & 77,982 & 4\\
24 & 140,683 & 140,691 & 8\\
25 & 254,487 & 254,504 & 17\\
26 & 461,409 & 461,439 & 30\\
27 & 838,433 & 838,487 & 54\\
28 & 1,526,536 & 1,526,635 & 99\\
\bottomrule
\end{tabular}
\caption{Exact finite certificates. All terms from $n=1$ through $n=21$ are
proved exactly by matching bounds. Completeness for $22\le n\le28$ is not
claimed. The full table is included as \texttt{data/term\_bounds.csv}.}
\label{tab:terms}
\end{table}

The unresolved gaps in this table do not weaken the exponential lower bound:
the digit certificate needs only distinct values with known expression costs.
Nor does the existence proof depend on trusting any particular OEIS term.

\section{Exact equality is decidable with an explicit bound}\label{sec:separation}

For completeness, we give a finite, entirely explicit alternative to an exact
algebraic-number package. It is much too pessimistic for practical large
computations, but establishes effectivity without numerical assumptions.

\begin{proposition}[Uniform separation]\label{prop:separation}
If $x\ne y$ are represented by expressions of size at most $n$, then
\begin{equation}\label{eq:separation}
 |x-y|\ge
 (n+1)^{-(2^{2n-2}-1)}=:\delta_n.
\end{equation}
\end{proposition}
\begin{proof}
Adjoin all radical subexpression values from both expressions. The resulting
field $K$ has degree at most $2^{2n-2}$: each expression has at most $n-1$
radical nodes, each requiring an extension of degree at most two.
The nonzero number $x-y$ is an algebraic integer, so
$|\Norm_{K/\Q}(x-y)|$ is a positive integer, at least one.
By Lemma~\ref{lem:basic}, under every embedding of $K$ its absolute value is
at most the sum of the two leaf counts, hence at most $n+1$.
One factor in the norm is the specified real value $x-y$. Bounding all the
other factors gives
$1\le|x-y|(n+1)^{[K:\Q]-1}$, which implies \eqref{eq:separation}.
\end{proof}

At precision $P$, fixed-denominator interval propagation as in
\eqref{eq:intervalsqrt} produces widths at most $2n\,2^{-P}$.
One can see this inductively: addition adds widths, and square root on
$[1,\infty)$ reduces a width by at least a factor of two before adding
less than two units of outward rounding.
Choose $P$ so that
\begin{equation}\label{eq:safeprecision}
 2^P>\frac{8n}{\delta_n}.
\end{equation}
Every interval then has width less than $\delta_n/4$.
If two such intervals overlap, their true values differ by less than
$\delta_n/2$, so Proposition~\ref{prop:separation} forces equality.
Disjoint intervals certify inequality. This gives a finite exact comparison
algorithm for all expressions of bounded size.

Alternatively, real algebraic numbers represented by defining polynomials and
isolating intervals provide exact equality and ordering; SageMath's
\texttt{AA} is a documented implementation \cite{sage}.
The present archive does not require SageMath. Its executed verifiers use only
the Python standard library, and the lower-bound verifier does not need the
astronomical worst-case precision \eqref{eq:safeprecision} because it never
declares overlapping intervals equal.

\section{Interpretation and remaining asymptotic questions}

The distinction between four levels of information is important.
First, exponential upper and lower estimates only show that
$\log a(n)=\Theta(n)$.
Second, the injective-spine argument proves the stronger assertion that the
coefficient of $n$ in $\log a(n)$ converges.
Third, the large-expression decomposition proves computability of that
coefficient and the exact critical identity \eqref{eq:criticalsum}.
Fourth, a full asymptotic equivalent would require further information that
is not supplied by these arguments.

In particular, the following issues remain unresolved in this article:
\begin{enumerate}[label=(\roman*),itemsep=3pt]
\item A substantially tighter numerical enclosure of $\gamma$.
The finite-data procedure is convergent, but its elementary upper codes are
not optimized for speed, and complete exact enumeration is expensive.
\item The existence of a ratio limit $a(n+1)/a(n)\to\gamma$.
A root limit and monotonicity alone do not imply this.
\item An asymptotic equivalent, including a power of $n$ and a leading
constant. The familiar $n^{-3/2}$ behavior of many tree models motivates a
question, not a conclusion, for the quotient by all numerical identities.
\end{enumerate}

It would be incorrect to identify the growth constant of an unordered-tree
majorant with $\gamma$. The majorants deliberately keep numerically redundant
expressions, and even rare local identities can change an exponential growth
constant. Likewise, a fit to the first 28 numbers cannot certify a limiting
base or a polynomial prefactor.

The rigorous outcome is
\[
 \boxed{\log a(n)=n\log\gamma+o(n),\quad
 1.86204<\gamma<1.9208,\quad
 \sum_{n\ge1}a(n)\gamma^{-n}=\gamma^2.}
\]
The constant is characterized by convergent algebraic brackets computed from
finite exact data. The finite numerical work in the archive checks the stated
bounds independently of approximate equality tests.

\appendix
\section{Reproducing the computations}

The archive is self-contained apart from a Python 3.10 or later installation and a LaTeX
installation for rebuilding the PDF. No computer-algebra package, internet
connection, or floating-point equality oracle is needed for the executed
certificates.

From the archive's top-level directory, run:
\begin{lstlisting}
python verify_certificate.py --output data/verification_results.json
python certify_upper_bound.py --output data/upper_bound_certificate.json
python analyze_counts.py
\end{lstlisting}
The first command replays all expression nodes, checks symbol costs, and
verifies pairwise disjoint enclosures. The second recomputes the formal upper
coefficients, the strict convergence inequality, and the improved finite-term
upper bounds. The third reproduces the rational lower-root brackets and the
CSV table.

To regenerate the lower-bound sample from scratch:
\begin{lstlisting}
python certified_enumerate.py --max-n 28 --precision 256 --output data
\end{lstlisting}
This generator is deliberately conservative about overlap. Its mathematical
specification is a certified sample, not a complete exact enumeration.
The independently replayable certificate is the primary finite witness.

In the execution used for this article, the sample generator completed in
about five seconds and the independent sample verifier in about two and a
half seconds. These are descriptive timings for that environment, not
portable performance guarantees. The certificate itself has 1,526,536
records, nine bytes per record, plus a 16-byte header.

\subsection*{Files in the archive}

{\small
\begin{longtable}{@{}p{0.42\textwidth}p{0.53\textwidth}@{}}
\toprule
File & Purpose\\
\midrule
\endhead
\texttt{article.tex}, \texttt{article.pdf} & This article and its source.\\
\texttt{README.md}, \texttt{Makefile} & Reproduction instructions and PDF build target.\\
\texttt{certified\_enumerate.py} & Conservative interval-based construction of the sample.\\
\texttt{verify\_certificate.py} & Independent integer-arithmetic replay and distinctness check.\\
\texttt{certify\_upper\_bound.py} & Exact coefficient recurrences and outward-rounded upper certificate.\\
\texttt{analyze\_counts.py} & Rational root brackets and the finite-term bounds table.\\
\texttt{data/digit\_sample.bin} & Compressed-by-sharing expression DAG; formula costs are charged in full.\\
\texttt{data/certificate\_metadata.json} & Generator counts and certificate interpretation.\\
\texttt{data/verification\_results.json} & Independently verified counts, minimum gap, and lower inequality.\\
\texttt{data/upper\_bound\_certificate.json} & Exact signed convergence witness and finite upper counts.\\
\texttt{data/finite\_lower\_roots.json} & Rational brackets for Table~\ref{tab:lower}.\\
\texttt{data/term\_bounds.csv} & All 28 lower and upper finite-term bounds.\\
\bottomrule
\end{longtable}
}

\section{Certificate format and verification boundaries}

The binary file begins with the eight ASCII bytes \texttt{A158415D}, followed
by two little-endian unsigned 32-bit integers: interval precision and maximum
assigned symbol cost. Each subsequent record consists of one unsigned byte
for cost and two unsigned 32-bit child indices.
The value \texttt{0xffffffff} is a sentinel. The first record has cost one
and two sentinels, representing $1$. A later record whose right index is the
sentinel represents the square root of its left child; otherwise it represents
the sum of its two children. Child indices must precede the current record.
The verifier recomputes each cost and enclosure rather than trusting them.

These checks establish concrete finite claims. The growth-limit,
critical-identity, and computability theorems are justified by the mathematical
proofs in the article; the archive does not contain a proof-assistant
formalization of them. Conversely, the numerical bounds do not depend on
assuming that the theoretical seed constructions have been explicitly
materialized as short expressions.

\begin{thebibliography}{9}

\bibitem{oeis}
V. Reshetnikov and contributors,
\emph{A158415: distinct values of expressions formed from $1$, square root,
and addition}, The On-Line Encyclopedia of Integer Sequences.
Entry introduced March 18, 2009; accessed September 20, 2026.
\url{https://oeis.org/A158415}.

\bibitem{discussion}
R. Kelly, R. Gerbicz, V. Reshetnikov, and others,
\emph{Comments on A158415}, archived sequence discussion, March 2009.
\url{https://oeis.org/A158415/a158415.txt}.

\bibitem{milne}
J. S. Milne,
\emph{Algebraic Number Theory}, version 3.08, July 19, 2020.
See especially Chapters 2, 3, and 7, including the discussion of Newton
polygons on printed page 125.
\url{https://www.jmilne.org/math/CourseNotes/ANT.pdf}.

\bibitem{fs}
P. Flajolet and R. Sedgewick,
\emph{Analytic Combinatorics}, Cambridge University Press, 2009.
Author-maintained booksite and full-text access:
\url{https://ac.cs.princeton.edu/home/}.

\bibitem{sage}
The Sage Development Team,
\emph{Algebraic Numbers and Number Fields}, reference documentation for
\texttt{AA} and \texttt{QQbar}; accessed September 20, 2026.
\url{https://doc.sagemath.org/html/en/reference/number_fields/sage/rings/qqbar.html}.

\end{thebibliography}
\end{document}
